\cotefile{157-4}
\scanpage{34}\mspage{27}
% archive sheet number: 33
3. Catégories de Quillen (deuxième mouture)

On se donne toujours un couple
\[
(M,W),\quad W\subset\operatorname{Fl}M,\qquad W\ \text{fblt saturé}\ \deleted{\text{\unsure{stable}}}.
\]
\inserted{$M$ satisfait $\unsure{\mathrm{M03}|\mathrm{MO3}|\mathrm{M}0_{3}}$ (stabilité par $\unsure{\varinjlim|\lim}$ \unsure{quelconques|qqconques} et $\unsure{\varprojlim|\lim}$ finies)}

\uline{Définition 1}\unsure{:} Une \uline{structure de Quillen} \uline{\unsure{faible|\deleted{faible}}} associée à la cat. de modèles $(M,W)$ est \deleted{\unsure{formée de}} \inserted{la donnée} de quatre parties $(C,TC,F,TF)$ de $\operatorname{Fl}(M)$, satisfaisant les conditions suivantes
\[
\begin{aligned}
\mathrm{M}\,1\qquad & C\longleftrightarrow TF,\quad TC\longleftrightarrow F\\
&\text{(i.e. }TF\subset C_{*},\quad F\subset(TC)_{*}\text{),}
\end{aligned}
\]

$\mathrm{M}\,2$\qquad Propriété de factorisation de Quillen pour $(C,TF)$ et $(TC,F)$

$\mathrm{M}\,3$\qquad Stabilité de $C$ par \unsure{composition} et par cochangement de base, de $F$ par composition et chgt. de base -- $C$ et $F$ contiennent les iso

$\mathrm{M}\,4$\qquad Stabilité similaire pour $TC$ et $TF$
\[
\begin{aligned}
\mathrm{M}'\deleted{\unsure{5}}\qquad & TC\subset C\cap W,\qquad TF\subset F\cap W\\
\mathrm{M}''\qquad & C\subset\operatorname{Cof}_{W},\qquad F\subset\operatorname{Fib}_{W}
\end{aligned}
\]

On dit que la structure de Quillen est \uline{exacte} si elle satisfait à l'axiome suppl.
\[
\mathrm{M}\,2\ \mathrm{bis}:\qquad W=TF\circ TC.
\]
% slanted
\margin{\uline{NB} Si $C,F,W$ stables par facteurs directs, (\unsure{donc|comme} aussi $C\cap W$ et $F\cap W$) alors l'\unsure{exactitude} à gauche $TC=C\cap W$ \uline{équivaut} à l'exactitude à droite $TF=F\cap W$, et \unsure{implique|impliquent} \deleted{\unsure{aussi}} que $C,\unsure{TC},\unsure{TF|F},TF$ \unsure{sont} \uline{clos} (\unsure{à g} pour $C,TC$ \unsure{à dr pour} $\unsure{F,TF}$).}

On dit qu'elle est \uline{exacte à gauche} si elle satisfait l'axiome suppl.
\[
\unsure{\mathrm M'_0(g)}:\qquad TC=C\cap W\qquad \text{(au lieu de }TC\subset C\cap W\text{ seulement)}
\]
\uline{exacte à droite} si elle satisfait
\[
\mathrm M'_0(\unsure{d}):\qquad TF=F\cap W\qquad \text{(au lieu de }TF\subset F\cap W\text{ seulement)}
\]
\scanpage{35}\mspage{28}
% archive sheet number: 34
\uline{Remarques} a) Une structure de Quillen exacte à gauche ou exacte à droite est exacte.

b) Si on a une \deleted{structure de} donnée $(C,TC,F,TF)$ satisfaisant $\mathrm M1,\mathrm M2,\mathrm M',\mathrm M''$, quand on sature $C,TC$ à gauche et $F,TF$ à droite, pour les stabilités précisées dans $\unsure{\mathrm M4|\mathrm M1},\mathrm M3$, on trouve une structure de Quillen. (Tous les axiomes sont évidents, sauf $\mathrm M'$, pour lequel je renvoie à p. 9, 10, \unsure{notamment} prop. 4.) Si on avait $\mathrm M2$ bis, \unsure{cel|cela} sera encore vérifié (a fortiori) après saturation, i.e. on a une structure de Quillen exacte.

\uline{Définition 1 bis} La structure de Quillen faible est dite \uline{structure de Quillen} (tout court) si elle est exacte à gauche et à droite, ou (comme on dit) ``biexacte''\unsure{;} \inserted{alors $(C,F)$ déterminent déjà $(C,TC,F,TF)$} (Ça correspond exactement aux ``model categories'' de Quillen.) On dit que la structure de Quillen faible est \uline{close} si on a
\[
\begin{aligned}
&C\mathrel{\leftrightarrow\mkern-15mu\mid\mkern4mu}TF,\qquad TC\mathrel{\leftrightarrow\mkern-15mu\mid\mkern4mu}F,\\
&\deleted{\unsure{TC}}TF=C_{*},\ C=(TF)^{*}\qquad\text{et}\quad F=(TC)_{*},\ TC=F^{*}.
\end{aligned}
\]
% slanted
\margin{Cela signifie aussi que $C,TF,TC,F$ sont stables par facteurs directs. On dit que la str. $\unsure{Qf|Q}$ est close à g. (resp. à dr.) si $C,TC$ (resp. $F,TF$) \unsure{sont} clos à g. (resp. à dr.)}

Une structure de Quillen faible à la fois biexacte (i.e. pas faible du tout !) et close est appelée str. de Qu. \uline{parfaite}. (Ça correspond exactement \unsure{à} \deleted{\unsure{aux|au}} ``closed model category'' de Quillen)
\scanpage{36}\mspage{29}
% archive sheet number: 35
\uline{Remarque} Pour qu'une structure de Quillen \inserted{faible} soit parfaite, il suffit qu'elle soit biexacte et qu'on ait
\[
 C=(TF)^{*},\qquad F=(TC)_{*}
\]
\deleted{$\unsure{TC=C\cap W}$} (ce qui implique déjà que $TC\unsure{\ne|=}F^{*},TF=C_{*}$, i.e. que $TC$ est close à g et $TF$ close à dr.) C'est prouvé dans Quillen. Ça implique aussi que $W$ est fortement saturé -- c'est le pays de cocagne !

\uline{Scholie}\unsure{:} Le couple $TC\longleftrightarrow F$, satisfaisant à la condition de factorisation de Quillen, et satisfaisant $TC\subset W,\ F\subset\operatorname{Fib}_{W}$, permet déjà de développer une théorie satisfaisante des fibres $W$-homotopiques dans $M$ -- mais sans \deleted{l'axiome} \inserted{propriété} d'exactitude pour les carrés $W$-cartésiens. Dualement, la donnée de $C\unsure{\leftarrow|\leftrightarrow}TF$, avec factorisation, et $TF\subset W,C\subset\operatorname{\unsure{Cofibr|Cofib}}_{W}$, donne une bonne théorie des cofibres $W$-homotopiques, mais sans \deleted{axiome} \inserted{propriété} d'exactitude pour les carrés $W$-cocartésiens. La structure de Quillen faible réunit ces deux données duales l'une de l'autre, avec en plus le lien \deleted{suppl} \inserted{supplémentaire} entre ces deux données, par $TC\subset C,TF\subset F$ (stipulé dans $\mathrm M'$). \deleted{La donnée \unsure{du côté|du seul} $TC\leftrightarrow F$} \inserted{Cette donnée assure} déjà que
\[
\underbrace{M_{cf}}_{\substack{\text{sous-catégorie pleine de }M\text{ formée des objets fibrants et cofibrants}\\\text{à la fois}}}\,W_{cf}^{-1}\xrightarrow{\approx}MW^{-1}=\operatorname{Ho}_{W}
\]
\scanpage{37}\mspage{30}
% archive sheet number: 36
est une équivalence, mais non pas que le foncteur localisation
\[
 M_{cf}\unsure{-|\longrightarrow}M_{cf}W_{cf}^{-1}
\]
est un \deleted{\unsure{simple}} passage \unsure{aux|au} quotient sur les flèches, i.e. que si $X,Y$ sont dans $M_{cf}$, alors
\diagram{\begin{tikzcd}[ampersand replacement=\&,row sep=small]
\text{(a)}\qquad\operatorname{Hom}_{M}(X,Y)\arrow[r] \& \operatorname{Hom}_{W}(X,Y)\arrow[d,equal,"\mathrm{def}"]\\
\&\operatorname{Hom}_{\operatorname{Ho}_{W}}(X,Y)
\end{tikzcd}}
soit surjectif. Ceci est assuré par l'axiome d'exactitude $\mathrm M2$ bis
\[
 W=TF\circ TC.
\]

\deleted{\unsure{Ainsi}} Avec ce seul axiome, on n'a pas de caractérisation sympa de la relation d'équivalence dans $\operatorname{Hom}_{M}(X,Y)$ définie par la flèche (a) (pour $X,Y$ dans $M_{cf}$).

% slanted
\margin{Il semblerait que la relation de Q-homotopie à droite \deleted{\unsure{est}} est \deleted{\unsure{une}} \deleted{\unsure{équiv.}} une relation d'équivalence, pour $\unsure{X,Y}$ de $\unsure{M_{cf}|M_{f}}$, et cela pour la \unsure{seule}  str. de Quillen faible exacte à droite.}
Si on a exactitude à gauche
\[
 TC=C\cap W\qquad(\Longleftrightarrow\quad TC=C\cap W^{\unsure{\mathrm{univ}|\mathrm{uni}}}),
\]
alors \unsure{cette} relation d'équivalence est \deleted{\unsure{la}} engendrée \inserted{par la relation de} \uline{Q-homotopie à droite} (\unsure{impliquée} par la Q-homotopie à gauche), et la catégorie quotient
% slanted
\margin{\uline{NB} D'après \unsure{Baues|Barnes}, il semblerait que sur $M_{\unsure{cf|f}}$, cette relation d'homotopie à droite équivaut à la relation d'homotopie à gauche et que celle-ci \unsure{ne} dépend pas du ``cylindre'' choisi.}
\[
 \pi_{d}M_{cf}=\pi M_{cf}\approx\operatorname{Ho}_{W}\deleted{\unsure{M_{cf}}}
\]
s'envoie dans $\pi_{d}M_{c}$
\[
 \deleted{\unsure{\pi_{d}M_{c}}}\ \pi M_{cf}\longrightarrow\pi_{d}M_{c}
\]
et ce foncteur peut être interprété \unsure{comme} (\unsure{ce dernier|ceci}
\scanpage{38}\mspage{31}
% archive sheet number: 37
comme un \uline{adjoint à droite} (pl. fidèle) du foncteur de localisation
\[
 \pi_{d}M_{c}\longrightarrow(\pi_{d}M_{c})W_{c}^{-1}\approx\operatorname{Ho}_{W},
\]
et $\operatorname{Ho}_{W}$ se déduit de $\pi_{d}M_{c}$ par un \uline{calcul de fractions} \uline{à gauche}. Dualement, si on a exactitude à droite, \deleted{alors} i.e.
\[
 TF=F\cap W\quad(\Longleftrightarrow\quad TF=F\cap W^{\deleted{\unsure{\mathrm{univ}|\mathrm{uni}}}}_{\mathrm{univ}}),
\]
alors la relation d'équivalence pertinente sur $\operatorname{Hom}_{M}(X,Y)$, pour $X,Y$ dans $M_{cf}$, est la Q-\uline{homotopie à gauche} (impliquée par la Q-homotopie à droite), et cette fois la flèche \unsure{naturelle} de $\pi_{g}M_{cf}=\pi_{\deleted{\unsure{g}}}M_{cf}\approx\operatorname{Ho}_{W}$ dans $\pi_{d}M_{f}$
\diagram{\begin{tikzcd}[ampersand replacement=\&,row sep=small]
\pi M_{cf}\arrow[r] \& \pi_{d}M_f\\
\operatorname{Ho}_{W}\arrow[u,phantom,"\approx"{description,sloped}] \&
\end{tikzcd}}
est un \uline{adjoint à gauche} \inserted{(pl. fidèle)} du foncteur de localisation
\[
 \pi_{\unsure{d}}M_{f}\longrightarrow M_{f}W_{f}^{-1}\approx\operatorname{Ho}_{W},
\]
et \deleted{\unsure{la}}\unsure{sa|la} \unsure{localisée|localisation} s'explicite par un calcul de fractions à droite. \deleted{Dans le}

Dans le cas biexact (i.e. \inserted{celui} des véritables structures de Quillen), les deux relations de Q-homotopie dans $M_{cf}$ (\deleted{et} plus génér. dans \unsure{un} $\operatorname{Hom}_{M}(X,Y)$, si $X$ dans $M_{c}$ et $Y$ dans $M_{f}$)
\scanpage{39}\mspage{32}
% archive sheet number: 38
coïncident, et deviennent \inserted{$\unsure{=|\simeq}$} particulièrement simples. (Car si $f,g:X\unsure{=|\rightrightarrows}Y$, avec $X$ cofibrant et $Y$ fibrant, alors $f\sim g$ peut se tester sur \uline{n'importe} \uline{quel} ``objet cylindre'' pour $X$, ou sur \uline{n'importe quel} ``objet chemins'' pour $Y$.)

\inserted{Dans} le cas où la structure de Quillen faible est exacte à gauche, \deleted{et} i.e. où la relation d'équi\unsure{v.} pertinente dans $\operatorname{Fl}M_{cf}$ est la Q-homotopie à droite, alors l'axiome d'exactitude est vérifié pour les carrés $W$-cartésiens, mais je doute qu'il en soit de $=$ pour les carrés $W$-cocartésiens.$^{*}$
% slanted
\margin{Il semble bien \unsure{que c'est} le cas, d'après \unsure{Baues|Barnes}, \unsure{cf annotation|cf annotations} \unsure{marginale correspondante|marginales correspondantes} à page 30}
(Si on ne fait \deleted{\unsure{l'}} une hypothèse supplémentaire, p. ex. l'exactitude à dr. de la structure de Quillen) Si on n'a que l'exactitude (mais ni à g. ni à dr.), il semble qu'on ne peut en conclure l'exactitude ni des carrés $W$-cartésiens, ni des carrés $W$-cocartésiens, \deleted{On peut} \inserted{\unsure{donc}} pas celle des suites de \unsure{cosuspension|cosuspensions} (loop \unsure{sequence|sequences}), ni des suites de suspension. On pourrait donc penser que la condition d'exactitude, à elle seule, ne sert pas à grand-
\scanpage{40}\mspage{\unsure{33|35}}
% archive sheet number: 39
chose pour le formalisme homotopique dans les catégories de fractions $\operatorname{Ho}_{W}$ etc. Mais on a \inserted{pourtant} ceci : soit
\diagram{\begin{tikzcd}[ampersand replacement=\&,row sep=0pt,column sep=small]
\&\&1\\
\unsure{\Upsilon|\Psi}={} \& 0\arrow[ur]\arrow[dr]\&\\
\&\&2
\end{tikzcd}}
, et supposons que \deleted{$(M,W)$ \uline{et}} \inserted{$(M,W)$ \uuline{et}} $(M(\Upsilon),W(\Upsilon))$ \deleted{\unsure{admette}}admettent une structure de Quillen faible exacte, \deleted{\unsure{et} $(M,W)$} Alors dans $M$ les carrés $W$-\deleted{\unsure{co}}cartésiens \unsure{sont|\deleted{étant}} \inserted{$W$-}exacts, \inserted{d'où p. \unsure{ex. ds|ex. des|ex. dans} $M$ \unsure{des} suites \uline{exactes} \inserted{longues} de cosuspension (\deleted{loop} exact loop sequences).} \deleted{Dualement, si (\unsure{Il suffit}} (\unsure{cette} hypothèse sur $(M(\Upsilon),W(\Upsilon))$ sert uniquement à assurer que toute flèche de $\operatorname{Ho}_{W(\Upsilon)}=M(\Upsilon)W(\Upsilon)^{-1}$ est isomorphe à une flèche provenant d'une flèche de $M(\Upsilon)$ lui-même. \deleted{\unsure{)}} Mais c'est là une propriété apparemment très délicate, et je ne connais d'autre façon de l'obtenir que via l'existence d'une structure de Quillen faible exacte.)
% slanted
\margin{\deleted{il semble que je me suis \unsure{embrouillé}, et je doute que \unsure{le|ce} résultat \unsure{indiqué|ici} soit exact.}}

Dualement, si $(M,W)$ et $(M(\unsure{\Psi|\Phi}),W(\Psi))$ admettent des structures de Quillen faibles \unsure{exacte|exactes}, alors les carrés $W$-cocartésiens de $M$ sont $W$-exacts, et on trouve (si $M$ est pointée) des suites exactes \uline{longues} de suspension.
\scanpage{41}\mspage{34}
% archive sheet number: 40
Ceci montre que si, pour une cat. de modèles $(M,W)$ donnée, on s'intéresse à ``faire de l'algèbre homotopique'' \deleted{pour} tous les couples $(M(I),W(I))$ à la fois, avec $I$ dans $\operatorname{Cat}$, il suffit (si on veut avoir une bonne théorie des fibres \inserted{et cofibres} $W$-homotopiques dans ces catégories, y inclus la propriété de $W$-exactitude pour les carrés $W$-cartésiens, et pour les carrés $W$-cocartésiens, il suffit de disposer \deleted{pour} chaque $(M(I),W(I))$ d'une structure de Quillen faible exacte sans plus.

Ceci nous montre donc l'intérêt de développer des critères, pour une $(M,W)$, assurant qu'il en est bien ainsi. À ce sujet, je note que je ne sais \inserted{comment} vérifier l'exactitude d'une structure de Quillen \uline{que} \unsure{si|\deleted{si}} \unsure{je|\deleted{je}} dispose de l'exactitude à g, ou à dr. C'est dire que je \unsure{dois|sais} dégager\deleted{\unsure{pour}} en fait des conditions sur $(M,W)$ qui assurent (disons) que les $(M(I),W(I))$ admettent des structures de Quillen faibles exactes \uline{à gauche}.
\scanpage{42}\mspage{35}
% archive sheet number: 41
\begingroup\makeatletter\tagsleft@true\makeatother
Les structures de Quillen faibles \deleted{pour} les $(M(I),W(I))$ seront déduites d'une structure \deleted{\unsure{de} $\unsure{Qf}$ pour} $(M,W)$, de la façon suivante : elles seront toutes ``\uline{closes à droite}'' i.e.
\[
 \deleted{\unsure{M}}\ F(I)=(TC(I))_{*},\quad TF(I)=(C(I))_{*},\tag*{(\unsure{1})}
\]
\endgroup
en particulier (pour $I=e$)
\[
 F=(TC)_{*},\quad TF=C_{*}
\]
ce qui implique déjà (puisque $C\subset\operatorname{Cof}_{W},TC\subset W^{\unsure{\mathrm{univ}|\mathrm{uni}}}$)
\[
 (\operatorname{Cof}_{W})_{*}\subset W_{\deleted{\unsure{\mathrm{univ}}}},\quad(W^{\mathrm{univ}})_{*}\subset\operatorname{Fib}_{W},\tag*{(\unsure{2})}
\]
\inserted{(d'où $(\operatorname{Cof}_{W})_{*}\subset W_{\mathrm{univ}}$).}
En vertu de ça, pour définir les structures $Qf$ \deleted{pour les} $(M(I),W(I))$, il suffit de définir $C(I),TC(I)$, et $=$, comme on veut l'exactitude à gauche i.e.
\[
 TC(I)=C(I)\cap W(I)\tag*{(\unsure{3})}
\]
(en particulier
\[
 TC=C\cap W\qquad),\tag*{(\unsure{4})}
\]
il reste à définir $C$. On \deleted{\unsure{prendra}} \inserted{définira} $C(I)$ ``argument \deleted{\unsure{et}} par argument'', i.e. pour $f\in\operatorname{Fl}(M(I))$
\[
 f\in C(I)\Longleftrightarrow\alpha_{s}^{*}(f)\in C\quad\text{pour tt }s\in\operatorname{Ob}I,\tag*{(5)}
\]
où $\alpha_{s}:\deleted{\unsure{\{s\}}}e\longrightarrow I$ défini par $\alpha_{s}(e)=s$.
Il en résulte d'ailleurs aussitôt, \deleted{par} (3) \deleted{\unsure{et}}
\[
 f\in TC(I)\Longleftrightarrow\alpha_{s}^{*}(f)\in TC\quad\text{pour tt }s\in\operatorname{Ob}I.\tag*{(6)}
\]
\scanpage{43}\mspage{36}
% archive sheet number: 42
On suppose donc que $(M,W)$ admet une str. de $Qf$ \inserted{exacte à g. et} close à droite $(C,TC,F,TF)$, définie en termes de
\[
 C\subset\operatorname{Fl}M),\qquad\boxed{\deleted{\unsure{C\subset\operatorname{Cof}_{W}}}}
\]
par
\[
 TC=C\cap W,\quad F=(TC)_{*},\quad TF=C_{*}.
\]
Pour que cela donne bel et bien une str. de $Qf$ (forcément exacte à g et close à droite), il faut simplmt \deleted{vérifier} vérifier les conditions suivantes

\textcircled{\scriptsize i} Factorisation pour $(C,TF=C_{*})$ et $(TC=C\cap W,F=(TC)_{*})$ -- on reviendra sur ce point.

\textcircled{\scriptsize ii} Stabilité de $C$ par composition et cochgt. de base, et $C$ contient les iso (d'où résulte la stabilité similaire de $TC=C\cap W$, compte tenu que $C\subset\operatorname{Cof}_{W}$) (NB les stabilités relatives pour $F,TF$ sont triviales par définition $F=(TC)_{*},TF=C_{*}$)

\[
\begin{aligned}
\text{\textcircled{\scriptsize iii}}\quad&
\deleted{\unsure{(\operatorname{Cof}_{W})_{*}\subset\operatorname{Fib}_{W},\ (C\cap W)_{*}\subset C}}\quad C\subset\operatorname{Cof}_{W},\\
&\underset{\substack{\Vert\\\inserted{TF}}}{\overset{\Vert}{C_{*}}}\subset W,\qquad \deleted{\unsure{TF}}\ (C\cap W)_{*}\subset\operatorname{Fib}_{W}
\end{aligned}
\]

Pour vérifier, de $=$, \deleted{\unsure{des}} que \deleted{les} $(C(I),TC(I),F(I),TF(I))$ \deleted{\unsure{sont des}} \inserted{une} str. de $Qf$ pour \deleted{\unsure{le}} \deleted{$M(I)$} $W(I)$, il faut vérifier des conditions similaires : \textcircled{\scriptsize i} \textcircled{\scriptsize ii} \textcircled{\scriptsize iii}. \deleted{Pour} Or, (ii) pour $C(I)$ résulte aussitôt de la condition (ii) pour $C$. Il reste à examiner \unsure{(iii)} et (i).
\scanpage{44}\mspage{37}
% archive sheet number: 43
\begingroup\makeatletter\tagsleft@true\makeatother
Les inclusions de (iii).

L'inclusion
\[
 C\subset\operatorname{Cof}_{W}
\]
implique que toute $f\in C(I)$ est dans $\operatorname{Cof}_{W}$ argument par argument. Cela \unsure{indique|implique|veut que|veut dire} que $f$ est dans $\operatorname{Cof}_{W(I)}$, comme on voit aussitôt sur la définition $\leftarrow$ je l'explicite en le

\uline{Lemme 1} : Soit $f$ ou $X_{\unsure{\cdot}}\unsure{-|\to}Y_{\cdot}$ une flèche de $M(I)$. Pour que $f$ soit dans $\operatorname{Cof}_{W(I)}$ (resp. dans $\operatorname{Fib}_{W(I)}$), il \uline{suffit} que pour tt $s$ dans $\operatorname{Ob} I$, $\alpha_{s}^{*}(f)$ soit dans $\operatorname{Cof}_{W}$ (resp. dans $\operatorname{Fib}_{W}$). Idem (pour mémoire) pour \unsure{la} condition pour que \deleted{\unsure{\(f\)}} $f\in W(I)$ -- là il \uline{f et s} que $\alpha_{s}^{*}(f)\in W\quad\forall s\in\operatorname{Ob}I$.

Nous allons appliquer ce $=$ lemme pour \deleted{\unsure{ét}} établir les inclusions
\[
 C(I)_{*}\subset W(I),\quad (C(I)\cap W(I))_{*}\subset\operatorname{Fib}_{\deleted{\unsure{I}}W}(I)
\]
Donc il suffit de voir que
\[
\begin{aligned}
 &\alpha_{s}^{*}(C(I)_{*})\subset C_{*}\quad(\deleted{\unsure{\alpha_s}}\subset W)\\
 &\deleted{\unsure{\alpha_{s}^{*}((C(I)\cap W)_{*}))\subset}}\\
 &\alpha_{s}^{*}((\underbrace{C(I)\cap W(I)}_{(C\cap W)(I)})_{*})\subset(C\cap W)_{*}\quad(\subset\operatorname{Fib}_{W}),
\end{aligned}
\]
et plus généralement
\[
 \alpha_{s}^{*}(\unsure{\Phi|\mathscr P}(I)_{*})\subset\Phi_{*}\tag*{(7)}
\]
pour tte \inserted{(!)} partie $\Phi$ de $\operatorname{Fl}(M)$, définissant donc $\Phi(I)\subset\operatorname{Fl}(M(I))$ terme à terme.
\endgroup
\scanpage{45}\mspage{38}
% archive sheet number: 44
En explicitant la chose, on trouve la condition
\[
 \alpha_{s!}(\unsure{\Phi|\mathscr P})\subset\Phi(I)\qquad\text{pour tt }s\in\operatorname{Ob} I,
\]
pour que ça marche, \deleted{Mais i.e.,} on veut que
\[
\begin{aligned}
 \deleted{\unsure{f}}i\in\Phi&\Longrightarrow\alpha_{s!}(i)\in\Phi(I)\qquad\text{i.e}\\
 &\alpha_{t}^{*}(\alpha_{s!}(i))\in\Phi\qquad\forall t\in\operatorname{Ob}I.
\end{aligned}
\]
Or le foncteur $\alpha_{t}^{*}\alpha_{s!}$ s'explicite comme
\[
 \unsure{A|A'}\longmapsto\underbrace{\operatorname{Hom}_{I}(t,s)\times A}_{\substack{\text{somme de }\operatorname{Hom}_{I}(t,s)\\\text{copies de }A\text{ dans }M}},
\]
donc il faut supposer $\Phi$ \uline{stable} \uline{par} sommes directes (\unsure{du type envisagé|des types envisagés} $\operatorname{Hom}_{I}(t,s)\times A$). C'est bien le cas pour \inserted{$\Phi=C$ et $\Phi=C\cap W$, et pour} des sommes directes \uline{finies}, grâce à (ii).

\deleted{\illegible{1}} On trouve donc le

\uline{Lemme 2}\unsure{,} a) Soit $I$ dans $\operatorname{Cat}$ telle que les sommes directes du type $\operatorname{Hom}_{I}(t,s)$ existent dans $M$. (P. ex. on suppose que $I$ est \uline{\unsure{finie|fini}}, plus gén. que les $\operatorname{Hom}_{I}(t,s)$ \deleted{\unsure{sont} $I$} \deleted{\unsure{sont}} sont des ens finis ; ou que les ``petites'' sommes directes existent dans $M$.) Soit $\Phi$ une partie de $\operatorname{Fl}(M)$ stable par les types de sommes directes, \deleted{Alors} \inserted{\unsure{et soit|et}} $\Phi(I)\subset\operatorname{Fl}(M(I))$ la partie correspondante. Alors pour tt $s\in\operatorname{Ob} I$ on a
\scanpage{46}\mspage{39}
% archive sheet number: 45
\[
 \alpha_{s}^{*}(\unsure{\Phi|\mathscr P}(I)_{*})\subset\Phi_{*}.
\]
b) Supposons que les $\operatorname{Hom}_{I}(t,s)$ soient finis (p. ex. $I$ finie). Alors \deleted{\unsure{les}} l'hyp. de a) s'applique pour $\Phi\unsure{=|\in}C$ et pour $\Phi=C\cap W$, \deleted{donc} \deleted{\unsure{et la structure}} la \deleted{structure} condition (iii) p \unsure{35|36} est satisfaite pour la structure $(\inserted{W(I),}C(I),TC(I),F(I),TF(I))$ dans $M(I)$. Même conclusion sans hyp. de finitude sur les \deleted{$M(I)$,} \inserted{$\operatorname{Hom}_{I}(t,s)$,} quand on suppose $M$ stable par ``petites'' sommes directes, et $C$ et $\deleted{\unsure{TC}}W$ stables également par de telles sommes directes.

\uline{NB} Vu le \unsure{Scholie|scholie} plus haut, il est parfaitement \unsure{raisonnable|vraisemblable} de se borner \unsure{au besoin|en théorie} aux catégories $I$ \uline{finies} -- vu que \deleted{\unsure{les}} \unsure{catégorie} de ces $I$ est stable par les produits $I\unsure{\mapsto|\vdash}I\times\unsure{\Psi|\Upsilon}$, $I\unsure{\mapsto|\vdash}I\times\Phi$.

Il reste, pour être heureux, à vérifier la condition de factorisation \textcircled{\scriptsize i} pour $(C,TF),(TC,F)$ certes, mais aussi et, surtout pour les $(C(I),TF(I)),(TC(I),F(I))$. C'est dire que de partir \deleted{\unsure{d'une}} d'une structure de Quillen close (pays de cocagne)
\scanpage{47}\mspage{40}
% archive sheet number: 46
ne nous avance guère, \inserted{à lui \deleted{tout} seul} faute de savoir que l'axiome de factorisation se conserve, en passant de $(M,W\ldots)$ à $(M(I),W(I)\ldots)$.

C'est ici que le th. de factorisation du n° 1 \inserted{doit} s'avérer providentiel. Mais pour pouvoir l'appliquer aux couples $(C(I),C(I)_{*})$, $(TC(I),TC(I)_{*})$, il faut (en plus de l'accessibilité de $M$, ce qui est une condition fort modérée) des conditions de stabilité très fortes pour $C(I),TC(I)$.$^{*}$
% slanted
\margin{plus les \inserted{\uline{petit}s} syst. de générateurs (pour $C,TC$.}
À vrai dire, il \unsure{reste|peut} la stabilité de $C$ et de $C\cap W$ par limites inductives ordinales (ce qui implique la $=$ stabilité pour les $C(I)$, \deleted{$\unsure{C(I)\cap W}$} $(C\cap W)(I)=C(I)\cap W(I)$).

À dire vrai, je vois peu de chances de trouver une str. de $Qf$ $(C,TC,F,TF)$ associée à $W$, et satisfaisant \deleted{\unsure{aux}} les conditions dites y compris celles de stabilité de $C$ et de $C\cap W$ par $\varinjlim$ ordinales, sans supposer que $W$ elle-même ne soit déjà stable par \inserted{(petites)} limites inductives filtrantes.
\scanpage{48}\mspage{41}
% archive sheet number: 47
On peut songer à utiliser une factorisation \uline{fonctorielle} \unsure{des} flèches de $M$ en
\[
\begin{aligned}
f=pi\quad &\text{avec}\quad i\in TC,\ p\in F\\
&\text{ou}\quad i\in C,\ p\in TF,
\end{aligned}
\]
pour déduire une factorisation \unsure{simultanée} pour les flèches $\unsure{f.|f_{\cdot}}$ de $M(I)$ en
\[
\begin{aligned}
f=pi\qquad &i\in TC(I),\ p\in F(I)\\
&\text{ou}\quad i\in C(I),\ p\in TF(I),
\end{aligned}
\]
où $C(I),TC(I),F(I),TF(I)$ sont ici définis via $C,TC,F,TF$ ``argument par argument''. Mais l'ennui, c'est que en procédant de cette façon naïve, on n'aura plus nécessairement $C(I)\longleftrightarrow TF(I)$, $TC(I)\unsure{\longleftrightarrow|\longleftarrow} F(I)$. Il faut considérablement rapetisser $TF(I)$ et $F(I)$, pour avoir une propriété RLP \unsure{p.r.à|p.r.} $C(I)$ resp. $TC(I)$.

Donc je ne vois guère d'autre moyen de s'en tirer, dans ce contexte général, qu'en supposant la stabilité de $W$ par $\unsure{\varinjlim|\underline{\lim}}$ filtrantes (ce qui n'est après \unsure{tout} pas très méchant).
\scanpage{49}\mspage{42}
% archive sheet number: 48
Bon gré, mal gré, on revient à la situation envisagée dans le \S\,2. On va \unsure{délinéer|délimiter} l'essentiel de ce qui a été pressenti alors

\uline{Théorème \unsure{?|2}} (encore un peu conjectural) Soit $(M,W)$ une catégorie de modèles, avec $W$ (comme toujours) faiblement saturée, mais de plus on suppose $M$ stable par petites $\varinjlim$ et par $\varprojlim$ finies, et $W$ stable par $\varinjlim$ filtrantes, enfin $M$ accessible.

Considérons les conditions suivantes

a) Il existe une \unsure{structure|structur.} de $Qf$ pour $W$, exacte à gauche, close à droite, $(C,TC,F,TF)$.

a$'$) Idem a), mais avec \deleted{$\unsure{C,TC|C,F}$} de plus $C,TC$ clos à gauche (i.e. \inserted{on a une} structure de $Qf$ close), et \deleted{$C,TC$ \unsure{resp.}} $C,\inserted{TC}$ de la forme $\widetilde{C_{0}},\widetilde{TC_{0}}$, où $\unsure{C|C_{0}},TC_{0}$ sont petites.

a$''$) Il existe $C_{0}\subset\operatorname{Cof}_{W},\ TC_{0}\subset W^{\unsure{\mathrm{univ}|\mathrm{uni}}}$, parties petites, telles que $C_{0*}\subset W_{\deleted{\unsure{\mathrm{univ}}}},\ (TC_{0})_{*}\subset\operatorname{Fib}_{W}$.

b) \deleted{\unsure{C|les} \unsure{a}} \(\unsure{('|(1|l'}\)
\[
 (\operatorname{Cof}_{W})_{*}\subset W_{\mathrm{univ}},\quad(W^{\mathrm{univ}})_{*}\subset\operatorname{Fib}_{W}.
\]
(d'où résulte, $\operatorname{Cof}_{W}$ et $W^{\mathrm{univ}}$ étant clos à $\unsure{\theta|0|g}$, que $(W_{\mathrm{univ}})^{*}\subset\operatorname{Cof}_{W},\ (\operatorname{Fib}_{W})^{*}\subset W^{\mathrm{univ}}$).
